% Approfondissements Institut du 21 octobre 2026 ; aucune modification du calage historique.
\section{Unité de perte et incertitude jointe}
\label{sec:institut-unite-incertitude}

Les pertes historiques décrivent des incidents complets ; la construction additive en
tire une par pilier touché. Nous comparons quatre constructions, sans coûts observés
par pilier pour les départager, à fréquence, dispersion et canaux communs.

\subsection{Constructions de la perte d'un incident}
\label{sec:institut-unite-perte}
% SOURCES-SCRIPTS: 119 126

Notons $K$ le nombre de piliers touchés et $Y$ une perte nulle avec probabilité
$1-p_u$, égale à $u+\mathrm{GPD}(\xi,\sigma)$ sinon. La perte additive est
$\sum_{j=1}^K Y_j$, avec des $Y_j$ indépendantes. La perte moyenne divise cette somme
par $K$. L'incident unique ne tire qu'un $Y$ ; l'incident amplifié vaut $K^\eta Y$,
avec $\eta=1$ posé. La loi exacte de $K$ inclut le choc P4 actif. L'inversion de Fourier
donne les capitaux des seize configurations, leurs attributions de Shapley et l'interaction.

% Généré par actualiser_approfondissements_institut.py depuis 119.json.
% SOURCES-SCRIPTS: 119
\begin{center}\footnotesize
\renewcommand{\arraystretch}{1.15}
\begin{tabular}{@{}lrrrrr@{}}
\toprule
\textbf{Perte} & \textbf{Conforme} & \textbf{Non conforme} & \textbf{Écart} & \textbf{Rapport} & \textbf{Interaction} \\
\midrule
Additive & $5\,944$ & $19\,890$ & $13\,946$ & $3{,}35$ & $4\,410$ \\
Moyenne & $4\,468$ & $8\,718$ & $4\,250$ & $1{,}95$ & $292$ \\
Unique & $4\,774$ & $10\,720$ & $5\,946$ & $2{,}25$ & $976$ \\
Amplifiée & $6\,908$ & $27\,746$ & $20\,838$ & $4{,}02$ & $7\,148$ \\
\bottomrule
\end{tabular}
\end{center}
\noindent Montants en M\euro{} pour le périmètre financier étudié. Les quatre règles
utilisent le calage historique ; l'incident amplifié a un exposant posé égal à un.
Les déplacements sous raffinement de grille restent de l'ordre du M\euro{}.


La moyenne garde plusieurs Bernoulli et peut être inférieure à $u$. La propagation
et l'accumulation y contribuent négativement, en diluant une perte extrême.
Avec une perte unique, ces contributions sont nulles par construction : elles
modifient $K$, sans modifier $Y$. Cela ne réfute pas ces canaux ; leur effet
financier exige de mesurer le coût selon les piliers touchés.

L'interaction additive calculée sans simulation garde le signe de son estimation
Monte-Carlo. Ces comparaisons isolent l'unité de perte, sans remplacer les résultats historiques.

Les constructions additive et amplifiée diffèrent par la dépendance entre les $K$ composantes d'un
incident, à marges égales. Un mélange les relie : avec probabilité $\rho$, un choc commun les rend comonotones, c'est-à-dire déterminées par le même rang aléatoire dans leur loi marginale ; sinon elles sont indépendantes. De $\rho=0$ à $\rho=1$, le capital croît de $16\,\%$ à l'état conforme et
de $39\,\%$ à l'état non conforme. Le rapport entre états passe de $3{,}35$ à $4{,}02$ ($+20\,\%$), avec $3{,}74$
à $\rho=0{,}5$. La croissance est monotone sur $21$ valeurs de $\rho$ (constat numérique), et la perte
annuelle moyenne ne change pas.

Entre l'incident unique et la comonotonie, la dépendance porte $45\,\%$ de
l'écart de capital à l'état conforme et $46\,\%$ à l'état non conforme, le reste venant du nombre de
composantes. La marge de l'incident complet, réutilisée pour chaque composante, n'est pas testée.

\subsection{Bootstrap conjoint de calibration}
\label{sec:institut-bootstrap-joint}
% SOURCES-SCRIPTS: 120

À seuil fixé, nous rééchantillonnons les pertes entières de 2000--2026.
Chaque réplication recompte $p_u$, ajuste conjointement $\xi$ et $\sigma$, puis
applique ce triplet aux deux états et aux quatre constructions. Les percentiles
portent sur leurs écarts et rapports appariés, jamais sur une soustraction de bornes marginales.

% Généré depuis 120.json ; intervalles conditionnels, non incertitude totale.
% SOURCES-SCRIPTS: 120
L'échantillon compte $570$ pertes et $86$ excès. L'ajustement donne $\hat\xi=0{,}59523$, $\hat\sigma=57{,}98$~M\euro{} et $\hat p_u=0{,}150877$. Nous retenons $2\,000$ réplications, sans échec.

\begin{center}\footnotesize
\renewcommand{\arraystretch}{1.15}
\begin{tabular}{@{}lrrrr@{}}
\toprule
\textbf{Perte} & \multicolumn{2}{c}{\textbf{Écart (M\euro)}} & \multicolumn{2}{c}{\textbf{Rapport}} \\
 & \textbf{Ajustement} & \textbf{Bande à 95\,\%} & \textbf{Ajustement} & \textbf{Bande à 95\,\%} \\
\midrule
Additive & $13\,933$ & $[4\,872\,;90\,361]$ & $3{,}35$ & $[3{,}04\,;5{,}22]$ \\
Moyenne & $4\,248$ & $[1\,401\,;30\,606]$ & $1{,}95$ & $[1{,}77\,;2{,}97]$ \\
Unique & $5\,940$ & $[1\,782\,;33\,719]$ & $2{,}25$ & $[1{,}97\,;3{,}11]$ \\
Amplifiée & $20\,821$ & $[6\,718\,;100\,660]$ & $4{,}02$ & $[3{,}48\,;5{,}51]$ \\
\bottomrule
\end{tabular}
\end{center}
\noindent Bandes percentiles nominales à $95\,\%$, calculées directement sur les paires. La corrélation bootstrap entre $\xi$ et $\sigma$ vaut $-0{,}603$. Les $12$ ajustements avec $\xi\geq1$ sont conservés : leur moyenne est infinie, leur quantile reste défini. Le raffinement numérique des ajustements extrêmes et d'un ajustement médian ne déplace aucun capital de plus de $0{,}08\,\%$.


Ces bandes supposent les pertes échangeables, dont la loi jointe ne dépend pas de leur ordre, et le seuil fixé ; elles excluent sélection et
non-stationnarité, et laissent fixes la fréquence historique, la dispersion, les multiplicateurs et la
cascade d'expert. La dispersion des rapports montre que la forme de queue ne s'annule pas entre états.

\subsection{Mélange prédictif annuel}
\label{sec:institut-melange-predictif}
% SOURCES-SCRIPTS: 123

Pour chaque calibration du bootstrap, nous recalculons la loi annuelle complète aux deux
états, puis nous moyennons ces $2\,000$ lois. Le quantile du mélange est une construction
prédictive empirique, et non une loi prédictive bayésienne.

% Généré depuis 123.json ; mélange empirique, conditionnel aux hypothèses du moteur.
% SOURCES-SCRIPTS: 123
\begin{center}\footnotesize
\renewcommand{\arraystretch}{1.15}
\begin{tabular}{@{}lrrrrrr@{}}
\toprule
\textbf{Perte} & \multicolumn{2}{c}{\textbf{Conforme (M\euro)}} & \multicolumn{2}{c}{\textbf{Non conforme (M\euro)}} & \multicolumn{2}{c}{\textbf{Rapport}} \\
 & \textbf{Central} & \textbf{Mélange} & \textbf{Central} & \textbf{Mélange} & \textbf{Central} & \textbf{Mélange} \\
\midrule
Additive & $5\,940$ & $7\,030$ & $19\,872$ & $28\,991$ & $3{,}35$ & $4{,}12$ \\
Moyenne & $4\,463$ & $5\,217$ & $8\,712$ & $12\,314$ & $1{,}95$ & $2{,}36$ \\
Unique & $4\,770$ & $5\,478$ & $10\,710$ & $14\,040$ & $2{,}25$ & $2{,}56$ \\
Amplifiée & $6\,904$ & $7\,832$ & $27\,725$ & $35\,520$ & $4{,}02$ & $4{,}54$ \\
\bottomrule
\end{tabular}
\end{center}
\noindent Quantiles annuels à $99{,}5\,\%$, périmètre financier étudié. Le mélange moyenne les lois annuelles de $2\,000$ calibrations ; la lecture se fait sur une grille de résolution relative $0{,}26\,\%$.


\begin{figure}[htbp]
  \centering
  \figover{melange_predictif_annuel.png}
  \caption{Quantile annuel à $99{,}5\,\%$ sous la règle additive : fonction de répartition
  des quantiles conditionnels des $2\,000$ calibrations, calage central (pointillé) et
  quantile du mélange (trait plein).}
  \label{fig:melange-predictif}
\end{figure}

Sous la règle additive, le quantile du mélange dépasse le calage central de $18\,\%$ à
l'état conforme et de $46\,\%$ à l'état non conforme. L'écart grandit avec le niveau : à $99{,}9\,\%$,
l'état non conforme passe de $43\,734$ à $87\,360$. Les $12$ ajustements avec $\xi\geq1$ n'en
sont pas la cause : sans eux, le rapport du mélange vaut $4{,}05$. Le rapport $3{,}35$ est celui
de deux points de calage ; celui des quantiles du mélange vaut $4{,}12$. Aucun des deux ne
mesure un facteur vrai. Le mélange garde les hypothèses du bootstrap conjoint, et n'est pas un
backtest du quantile à $99{,}5\,\%$.

\subsection{Validation chronologique}
\label{sec:institut-validation-temporelle}
% SOURCES-SCRIPTS: 124

Sur les années complètes 2000--2025, soit $565$ pertes dont $84$ excès au-dessus du seuil fixé,
nous évaluons en origine glissante, de 2014 à 2025, trois lois de sévérité des excès : stationnaire,
à échelle variable et à deux régimes. Aucune loi variable n'améliore le score moyen de la stationnaire. L'écart de score logarithmique
moyen par excès vaut $-0{,}008$ pour l'échelle variable, d'intervalle bootstrap à $90\,\%$
$[-0{,}057\,;0{,}043]$, et $-0{,}025$ pour les deux régimes, $[-0{,}073\,;0{,}025]$.

Pour le taux de
dépassement du seuil, le gain moyen de la tendance logistique sur le taux constant est de $0{,}039$
par perte évaluée, $[0{,}001\,;0{,}076]$, avec une pente de $0{,}107$ par an.
Cet intervalle exploratoire rééchantillonne les pertes individuellement. Le contrôle par blocs
annuels du \S\ref{soc:sec:benchmark-apprentissage} donne un intervalle contenant zéro : l'incertitude sur le gain dépend du rééchantillonnage.

% Généré depuis 124.json ; validation chronologique, capital pour l'année suivante.
% SOURCES-SCRIPTS: 124
\begin{center}\footnotesize
\renewcommand{\arraystretch}{1.15}
\begin{tabular}{@{}lrrr@{}}
\toprule
\textbf{Loi retenue} & \textbf{Conforme} & \textbf{Non conforme} & \textbf{Rapport} \\
\midrule
Stationnaire & $5\,958$ & $20\,268$ & $3{,}40$ \\
Échelle variable & $6\,417$ & $21\,037$ & $3{,}28$ \\
Deux régimes & $5\,585$ & $18\,241$ & $3{,}27$ \\
Taux en tendance & $10\,782$ & $37\,647$ & $3{,}49$ \\
Échelle variable et taux en tendance & $11\,446$ & $38\,887$ & $3{,}40$ \\
\bottomrule
\end{tabular}
\end{center}
\noindent Capital annuel à $99{,}5\,\%$ pour 2026 (M\euro, règle additive, périmètre financier étudié), après ajustement sur 2000--2025.


Le niveau du capital varie de $-10\,\%$ à $+92\,\%$ selon la loi retenue, le rapport entre états de
$-4{,}0\,\%$ à $+2{,}6\,\%$. Le taux de dépassement prévu pour 2026 atteint $0{,}349$ contre $0{,}149$ :
extrapolation d'une tendance, sans validation. Les excès d'une même année sont traités comme indépendants, avec un à onze excès par
an : la puissance est faible, et ce n'est pas un backtest du quantile à $99{,}5\,\%$.

\subsection{Rééchantillonnage par années}
\label{sec:institut-reechantillonnage-annuel}
% SOURCES-SCRIPTS: 125

Nous tirons des années complètes de 2000 à 2025, avec leurs pertes, ou des blocs de trois années
consécutives, puis nous réestimons à chaque tirage le taux de dépassement et la queue. Le nombre
annuel de pertes enregistrées, de $4$ à $39$, module aussi la fréquence par
$\lambda_b=\lambda_{\mathrm{ref}}\,n_b/n_{\mathrm{ref}}$, ce qui conserve la référence historique.

% Généré depuis 125.json ; rééchantillonnage annuel conjoint, bandes percentiles exploratoires.
% SOURCES-SCRIPTS: 125
\begin{center}\footnotesize
\renewcommand{\arraystretch}{1.15}
\begin{tabular}{@{}llrrr@{}}
\toprule
\textbf{Tirage} & \textbf{Fréquence} & \textbf{Rapport médian} & \textbf{Bande à 90\,\%} & \textbf{Non conforme, bande à 90\,\%} \\
\midrule
Années & fixée & $3{,}36$ & $[3{,}01\,;4{,}23]$ & $[7\,259\,;51\,189]$ \\
Années & propagée & $3{,}36$ & $[3{,}00\,;4{,}22]$ & $[7\,192\,;52\,139]$ \\
Blocs de trois années & fixée & $3{,}19$ & $[2{,}95\,;4{,}06]$ & $[4\,710\,;46\,968]$ \\
Blocs de trois années & propagée & $3{,}19$ & $[2{,}96\,;4{,}07]$ & $[4\,644\,;50\,377]$ \\
\bottomrule
\end{tabular}
\end{center}
\noindent Règle additive, périmètre financier étudié, 600 tirages par schéma (capital en M\euro).


Propager la fréquence change peu. Le facteur $n_b/n_{\mathrm{ref}}$ reste entre $0{,}855$ et
$1{,}145$ à $5$ et $95\,\%$. La bande du rapport entre états est inchangée avec des années
échangeables ($+0\,\%$) et varie de $-1\,\%$ avec des blocs. Celle du capital non conforme
s'élargit de $2\,\%$, puis de $8\,\%$. Il s'agit de pertes enregistrées, non d'une fréquence par
entité ; la hausse du nombre annuel de pertes n'est pas modélisée.

\subsection{Précision numérique}
\label{sec:institut-fourier-controles}
% SOURCES-SCRIPTS: 119 120

L'inversion de Fourier calcule la loi de la somme annuelle à partir des transformées des pertes et du comptage. La grille rend les convolutions calculables par produits. L'inclinaison exponentielle atténue les masses lointaines, puis est retirée ; elle réduit leur repliement artificiel en début de grille. La masse hors portée
est omise, sans report dans la dernière case ni renormalisation. Le pas et la portée
s'adaptent aux ajustements extrêmes ; les contrôles les font varier avec l'inclinaison.
Pour $K=1$, les quatre règles coïncident. Une sévérité déterministe redonne la
binomiale négative amincie ; les moments tronqués restent compatibles avec les
moyennes analytiques finies.
